\section{From finite local models to long improving paths}\label{sec:lifting}
We first isolate the exact obstruction to the proposed stronger certificate tests. This matrix decomposition is classical, not a new theorem of this paper.

\begin{lemma}[PSD principal block; Shaked-Monderer \cite{SM}, Lemma 3.1]\label{lem:spn}
If $Z\in\COP_n$ and $Z_{[r],[r]}\succeq0$, then $Z\in\SPN_n:=\mathbb S_+^n+\mathbb N^n$, where $\mathbb N^n$ denotes the symmetric entrywise nonnegative matrices. Hence $\ip A Z\ge0$ for every doubly nonnegative $A$.
\end{lemma}
\begin{proof}
First let $Z=\left(\begin{smallmatrix}G&b\\b\trans&c\end{smallmatrix}\right)$ with $G\succ0$. Minimize $x\trans Gx+2b\trans x$ on $x\ge0$, with minimizer $z$ and residual $\ell=Gz+b\ge0$, $z\trans\ell=0$. Copositivity at $(z,1)$ implies $c\ge z\trans Gz$. The decomposition
\[
 Z=\begin{pmatrix}G&-Gz\\-z\trans G&z\trans Gz\end{pmatrix}
   +\begin{pmatrix}0&\ell\\\ell\trans&c-z\trans Gz\end{pmatrix}
\]
is PSD plus nonnegative. For singular $G$, apply this result to $Z+\varepsilon\diag(I_r,0)$ and let $\varepsilon\downarrow0$. The SPN cone is closed: in a convergent sequence $Z_j=P_j+N_j$, diagonal entries of $P_j\succeq0$ and $N_j\ge0$ are bounded; $|(P_j)_{ab}|\le\sqrt{(P_j)_{aa}(P_j)_{bb}}$ bounds $P_j$, and then $N_j=Z_j-P_j$ is bounded. A convergent subsequence gives a limiting decomposition. The pairing assertion follows from PSD and entrywise nonnegative duality.
\end{proof}

\begin{proof}[Proof of Theorem~\ref{thm:main}]
Write $A=(B_N+Z)/(2\beta N^2)$, where $B_N=S_NBS_N$ and $\norm Z_F\le\eta$. Since $N\ge1$, $B_N\succeq\mu I$ and every entry is at least $b_*$. The choice of $\eta$ preserves positive definiteness and strict entrywise positivity. Also $\norm A_F\le(\beta N^2+1)/(2\beta N^2)\le1$.

Positive diagonal congruence preserves copositivity. Thus $W_N$ has cone margin $t_B/N^2$, norm at most $w_B$, and
\[
 \ip{W_N}{B_N}\le-\gamma+\frac{t_B}{N^2}\tr(B_N)
 \le-3\gamma/4.
\]
The perturbation contributes at most $\eta w_B\le\gamma/4$. Division by $2\beta N^2$ gives the stated pairing and distance. In particular every ball input is outside the ordinary CP cone.

Suppose a vertex is reachable in at most $m$ bounded edges within $\mathcal F_H$ from one of the allowed starts. Lemma~\ref{lem:gap} gives
\[
 \norm b_\infty\le R_0+mL_H,\qquad
 \norm k_\infty\le\norm z_\infty
 \le K_H(R_0+mL_H+1/2),
\]
so $v_k\le s_m$. For a primitive integral extreme ray $R$, define
\[
 h(R)=\ip{B_0}{R_{[r],[r]}}.
\]
All leading active constraints of $H$ remain active. Since $B_0$ is a positive combination of all their normals,
\begin{equation}\label{eq:exposure}
 h(R)\ge0,\qquad h(R)=0\ \Longleftrightarrow\ R_{[r],[r]}=0.
\end{equation}
The reverse implication in the equality uses that the active squares span $\mathbb S^r$. If $h(R)>0$, integrality gives $h(R)\ge1/q_B$. Applying the sharper block bounds \eqref{eq:raybounds},
\[
 \begin{split}
 \ip{B_N+Z}R
 &\ge\frac{N^2}{q_B}
       -2N\norm a_1F_Hs_m
       -(\alpha+\eta n)F_Hs_m^2\\
 &\ge\frac{N^2}{2q_B}>0.
 \end{split}
\]
For the last inequality, $N\ge8q_BF_HT_Bs_m$, $F_H,q_B\ge1$, $\eta\le1$, and $T_B\ge\norm a_1,\alpha+n,1$ bound the two subtracted fractions by $1/4$ and $1/64$ of $N^2/q_B$, respectively. Thus every improving extreme ray has zero leading block. The position-dependent quadratic growth of its last entry does not require a quadratic scale $N$: it is dominated by the $N^2$ leading exposure once $N$ is linear in the shear parameter.

There is an improving extreme ray at every departure vertex. Otherwise $A$ belongs to the minimizing normal cone, a finite cone generated by active nonnegative integer squares, and hence to $\CP_n$, a contradiction. The tangent cone is full dimensional and pointed by local polyhedrality and perfection. Its improving extreme ray has zero leading block and is therefore tangent to the face. It cannot be a recession direction: by Lemma~\ref{lem:gap}, those directions have nonnegative cross and bottom entries, so the strictly entrywise positive objective has positive cost on every nonzero such direction. A finite endpoint exists by closedness and local polyhedrality throughout the face; it is the adjacent vertex. Lemma~\ref{lem:spn} prevents a negative pairing at the current vertex, and also rules out every copositive certificate with PSD leading block. This proves the induction for all choices through $m$ pivots.

For the stated history certificate, its leading block is $(\sum_j\lambda_j)H+Q_{[r],[r]}\succeq0$, because every visited ray has zero leading block. Lemma~\ref{lem:spn} applies even if the procedure optimizes the correction or the history coefficients using unrestricted computation. Arbitrary nonnegative diagonal corrections are included; in fact arbitrary PSD corrections are included. The argument does not merely bound a particular correction formula.

All constants are fixed for a fixed block, seed, and starting region, and $N_m=\Theta(m+1)$. Entries of $A_{N_m}$ and $W_{N_m}$ have $O(\log(m+1))$ bits. The last diagonal of $A_{N_m}$ is a fixed positive rational multiple of $N_m^{-2}$, giving the matching lower encoding bound after reducing fractions. The distance estimate gives the same upper bound on logarithmic inverse distance. The local-cost claim is Corollary~\ref{cor:localcost}. Restarting within the starting region repeats an uncompleted prefix rather than shortening it; history certificates of the displayed form remain blocked across runs.
\end{proof}

\begin{remark}[An abstract exposure lemma]
The confinement argument is not peculiar to matrix cones. On a rational locally polyhedral set, let a face be exposed on its tangent cones by an integral functional $h$ with gap $1/q$ on primitive extreme rays outside the face. If local models yield transverse components bounded by $C(1+t)$ and $C(1+t)^2$ within graph distance $t$, the objective $N^2h+N\ell+a$ excludes every escaping improving extreme ray for $N\ge C'(1+t)$. The matrix theorem supplies verifiable hypotheses, the finite local models, and certificate exclusion for an important nonpolyhedral cone; it is not merely a formal restatement of this elementary domination argument.
\end{remark}
